\scanpage{298}
\mspage{22}
% archive sheet number: 297

besoin apparemment de $\mathrm{Loc}(6'a\ \text{\uuline{et}}\ b)$.) On en conclut $C\subset\mathrm{Cof}_{\unsure{W|w}}$ -- c'est merveilleux ! (Mais je ne sais pas prouver le lemme \ldots)

Cela légitime la démonstration faite en \uline{XII} loc. cit., et on trouve la

\uline{Proposition} 3.3. La condition $\mathrm{Loc}(6')$ (axiome des carrés cocartésiens, forme forte\inserted{a \& b}) \deleted{implique} \uline{équivaut} à $\mathrm{Loc}(6.a,b)$ (axiome des hom. de carrés cocartésiens)\inserted{et est \uline{impliquée} par $\mathrm{Loc}(6c)$.} \deleted{\unsure{forme} \illegible{2}}
\deleted{(qui implique (\unsure{déjà, th?} 18) \unsure{d'après} prop. 3.1.)}
\inserted{Si $W$ satisfait à $\mathrm{Loc}(1\,\mathrm{bis})$, i.e. stable par fact. dir. alors $\mathrm{Loc}(6')$ équivaut et à $\mathrm{Loc}(6abc)$ \deleted{et à $\mathrm{Loc}(6c)$} \inserted{$\mathrm{Loc}(6,c)$.}}
% slanted
\margin{\uline{\unsure{Commentaire|Corollaire}}}

3.\deleted{\unsure{3}}4. Ainsi, on trouve une théorie des cofibres $W$-homotopiques \uline{soit} sous l'hypothèse $\mathrm{Loc}(6')$, (équivalente à $\mathrm{Loc}(6)$\inserted{plus \unsure{compliquée|compliqué}}), soit sous l'hypothèse \deleted{plus} faible\inserted{et $+$ compliquée} $\mathrm{Loc}(6a)$.

Dans le premier cas, c'est une théorie \inserted{de nature} très particulière, puisque les carrés $W$-cocartésiens sont ``réversibles'', (i.e. on a \deleted{\unsure{implicat}} non seulement $i_1\in W\Longrightarrow j_1\in W$, mais aussi l'implication inverse).

Il reste donc la question si $\mathrm{Loc}(6'a)$ \deleted{est équivalente} suffit pour une théorie des cofibres \inserted{$W$-}homotopiques, ou ce qui revient sûrement au $=$, si ça implique déjà $\mathrm{Loc}(6a)$ (donc lui est équivalent), C'est \unsure{entièrement|ultérieurement} suspendu à la \uline{question} :
